% =====================================================================
% Figura: carga puntual frente a un plano conductor a tierra.
% (a) Hipótesis errónea: líneas radiales de q aislada; atravesarían el
%     conductor (tramos rojos) y no llegan perpendiculares a él.
% (b) Solución real: las líneas mueren perpendicularmente sobre el plano;
%     cada marca azul es el pie de una línea (densidad ~ |sigma_ind|).
% Líneas con igual flujo por línea (cos alpha equiespaciado), d = 1.
% =====================================================================
\begin{tikzpicture}
  \def\esc{1.45}
  \newcommand{\planoConductor}{%
    \shade[top color=conductor,bottom color=fondo] (-2.3*\esc,0) rectangle (2.3*\esc,-1.0*\esc);
    \draw[gristexto,line width=0.7pt] (-2.3*\esc,0) -- (2.3*\esc,0);}
  % ---------------- (a)
  \begin{scope}
    \fill[panel] (-2.3*\esc,0) rectangle (2.3*\esc,2.75*\esc);
    \planoConductor
    \begin{scope}
      \clip (-2.3*\esc,-1.0*\esc) rectangle (2.3*\esc,2.75*\esc);
      % Generado automáticamente por generar_figuras.py; no editar.
\draw[flinea, flecha=1.514cm] plot coordinates {(0.046,1.540) (1.308,3.987)};
\draw[flinea, flecha=1.514cm] plot coordinates {(-0.046,1.540) (-1.308,3.987)};
\draw[flinea, flecha=2.038cm] plot coordinates {(0.076,1.518) (2.837,3.987)};
\draw[flinea, flecha=2.038cm] plot coordinates {(-0.076,1.518) (-2.837,3.987)};
\draw[flinea, flecha=1.992cm] plot coordinates {(0.091,1.495) (3.335,3.105)};
\draw[flinea, flecha=1.992cm] plot coordinates {(-0.091,1.495) (-3.335,3.105)};
\draw[flinea, flecha=1.825cm] plot coordinates {(0.099,1.473) (3.335,2.210)};
\draw[flinea, flecha=1.825cm] plot coordinates {(-0.099,1.473) (-3.335,2.210)};
\draw[flinea, flecha=1.778cm] plot coordinates {(0.102,1.450) (3.335,1.450)};
\draw[flinea, flecha=1.778cm] plot coordinates {(-0.102,1.450) (-3.335,1.450)};
\draw[flinea, flecha=1.825cm] plot coordinates {(0.099,1.427) (3.335,0.690)};
\draw[flinea, flecha=1.825cm] plot coordinates {(-0.099,1.427) (-3.335,0.690)};
\draw[flinea, flecha=1.739cm] plot coordinates {(0.091,1.405) (2.923,0.000)};
\draw[flinea, prohibida] plot coordinates {(2.923,0.000) (3.335,-0.205)};
\draw[flinea, flecha=1.739cm] plot coordinates {(-0.091,1.405) (-2.923,0.000)};
\draw[flinea, prohibida] plot coordinates {(-2.923,0.000) (-3.335,-0.205)};
\draw[flinea, flecha=1.140cm] plot coordinates {(0.076,1.382) (1.621,0.000)};
\draw[flinea, prohibida] plot coordinates {(1.621,0.000) (3.080,-1.305)};
\draw[flinea, flecha=1.140cm] plot coordinates {(-0.076,1.382) (-1.621,0.000)};
\draw[flinea, prohibida] plot coordinates {(-1.621,0.000) (-3.080,-1.305)};
\draw[flinea, flecha=0.841cm] plot coordinates {(0.046,1.360) (0.747,0.000)};
\draw[flinea, prohibida] plot coordinates {(0.747,0.000) (1.420,-1.305)};
\draw[flinea, flecha=0.841cm] plot coordinates {(-0.046,1.360) (-0.747,0.000)};
\draw[flinea, prohibida] plot coordinates {(-0.747,0.000) (-1.420,-1.305)};

    \end{scope}
    \node[cargapos] at (0,\esc) {};
    \node[rotulo,anchor=south,fill=panel,inner sep=1pt,rounded corners=1pt] at (0,1.0*\esc+0.17) {$q$};
    \node[rotulogris,anchor=south east,fill=conductorosc,inner sep=2pt,rounded corners=1pt]
       at (2.25*\esc,-0.95*\esc) {¡atraviesa el conductor!};
    \draw[azulborde,line width=0.5pt] (-2.3*\esc,-1.0*\esc) rectangle (2.3*\esc,2.75*\esc);
    \node[rotulogris,anchor=south west] at (-2.3*\esc,2.77*\esc)
      {(a) Campo de $q$ aislada (hipótesis errónea)};
  \end{scope}
  % ---------------- (b)
  \begin{scope}[shift={(7.35,0)}]
    \fill[panel] (-2.3*\esc,0) rectangle (2.3*\esc,2.75*\esc);
    \planoConductor
    \begin{scope}
      \clip (-2.3*\esc,-1.0*\esc) rectangle (2.3*\esc,2.75*\esc);
      % Generado automáticamente por generar_figuras.py; no editar.
\draw[flinea, flecha=1.305cm] plot coordinates {(0.046,1.540) (0.073,1.592) (0.100,1.643) (0.127,1.695) (0.153,1.746) (0.180,1.797) (0.208,1.849) (0.235,1.900) (0.262,1.951) (0.290,2.002) (0.317,2.053) (0.345,2.104) (0.373,2.155) (0.401,2.205) (0.430,2.256) (0.458,2.306) (0.487,2.357) (0.516,2.407) (0.545,2.457) (0.575,2.507) (0.604,2.557) (0.634,2.607) (0.664,2.657) (0.694,2.706) (0.724,2.756) (0.755,2.805) (0.785,2.854) (0.816,2.903) (0.847,2.952) (0.879,3.001) (0.910,3.050) (0.942,3.098) (0.974,3.147) (1.006,3.195) (1.038,3.243) (1.071,3.291) (1.104,3.339) (1.136,3.387) (1.170,3.435) (1.203,3.482) (1.236,3.529) (1.270,3.577) (1.304,3.624) (1.338,3.671) (1.372,3.717) (1.407,3.764) (1.441,3.811) (1.476,3.857) (1.511,3.903) (1.547,3.949) (1.582,3.995)};
\draw[flinea, flecha=1.305cm] plot coordinates {(-0.046,1.540) (-0.073,1.592) (-0.100,1.643) (-0.127,1.695) (-0.153,1.746) (-0.180,1.797) (-0.208,1.849) (-0.235,1.900) (-0.262,1.951) (-0.290,2.002) (-0.317,2.053) (-0.345,2.104) (-0.373,2.155) (-0.401,2.205) (-0.430,2.256) (-0.458,2.306) (-0.487,2.357) (-0.516,2.407) (-0.545,2.457) (-0.575,2.507) (-0.604,2.557) (-0.634,2.607) (-0.664,2.657) (-0.694,2.706) (-0.724,2.756) (-0.755,2.805) (-0.785,2.854) (-0.816,2.903) (-0.847,2.952) (-0.879,3.001) (-0.910,3.050) (-0.942,3.098) (-0.974,3.147) (-1.006,3.195) (-1.038,3.243) (-1.071,3.291) (-1.104,3.339) (-1.136,3.387) (-1.170,3.435) (-1.203,3.482) (-1.236,3.529) (-1.270,3.577) (-1.304,3.624) (-1.338,3.671) (-1.372,3.717) (-1.407,3.764) (-1.441,3.811) (-1.476,3.857) (-1.511,3.903) (-1.547,3.949) (-1.582,3.995)};
\draw[flinea, flecha=1.696cm] plot coordinates {(0.076,1.518) (0.119,1.556) (0.162,1.595) (0.206,1.633) (0.249,1.671) (0.293,1.709) (0.337,1.747) (0.381,1.785) (0.425,1.822) (0.470,1.860) (0.515,1.897) (0.560,1.933) (0.605,1.970) (0.650,2.006) (0.696,2.042) (0.742,2.077) (0.788,2.112) (0.834,2.147) (0.881,2.181) (0.928,2.216) (0.975,2.249) (1.022,2.283) (1.070,2.316) (1.118,2.348) (1.166,2.380) (1.215,2.412) (1.264,2.444) (1.313,2.475) (1.362,2.505) (1.411,2.535) (1.461,2.565) (1.511,2.595) (1.562,2.623) (1.612,2.652) (1.663,2.680) (1.714,2.707) (1.765,2.734) (1.817,2.761) (1.869,2.787) (1.921,2.813) (1.973,2.838) (2.025,2.863) (2.078,2.887) (2.131,2.911) (2.184,2.934) (2.237,2.957) (2.291,2.979) (2.345,3.001) (2.398,3.022) (2.453,3.043) (2.507,3.063) (2.562,3.083) (2.616,3.102) (2.671,3.121) (2.726,3.139) (2.781,3.157) (2.837,3.174) (2.892,3.191) (2.948,3.207) (3.004,3.222) (3.060,3.237) (3.116,3.252) (3.173,3.266) (3.229,3.279) (3.286,3.292) (3.342,3.304)};
\draw[flinea, flecha=1.696cm] plot coordinates {(-0.076,1.518) (-0.119,1.556) (-0.162,1.595) (-0.206,1.633) (-0.249,1.671) (-0.293,1.709) (-0.337,1.747) (-0.381,1.785) (-0.425,1.822) (-0.470,1.860) (-0.515,1.897) (-0.560,1.933) (-0.605,1.970) (-0.650,2.006) (-0.696,2.042) (-0.742,2.077) (-0.788,2.112) (-0.834,2.147) (-0.881,2.181) (-0.928,2.216) (-0.975,2.249) (-1.022,2.283) (-1.070,2.316) (-1.118,2.348) (-1.166,2.380) (-1.215,2.412) (-1.264,2.444) (-1.313,2.475) (-1.362,2.505) (-1.411,2.535) (-1.461,2.565) (-1.511,2.595) (-1.562,2.623) (-1.612,2.652) (-1.663,2.680) (-1.714,2.707) (-1.765,2.734) (-1.817,2.761) (-1.869,2.787) (-1.921,2.813) (-1.973,2.838) (-2.025,2.863) (-2.078,2.887) (-2.131,2.911) (-2.184,2.934) (-2.237,2.957) (-2.291,2.979) (-2.345,3.001) (-2.398,3.022) (-2.453,3.043) (-2.507,3.063) (-2.562,3.083) (-2.616,3.102) (-2.671,3.121) (-2.726,3.139) (-2.781,3.157) (-2.837,3.174) (-2.892,3.191) (-2.948,3.207) (-3.004,3.222) (-3.060,3.237) (-3.116,3.252) (-3.173,3.266) (-3.229,3.279) (-3.286,3.292) (-3.342,3.304)};
\draw[flinea, flecha=1.520cm] plot coordinates {(0.091,1.495) (0.143,1.521) (0.195,1.546) (0.247,1.572) (0.299,1.597) (0.352,1.622) (0.404,1.646) (0.457,1.670) (0.510,1.694) (0.563,1.718) (0.616,1.741) (0.670,1.763) (0.723,1.786) (0.777,1.807) (0.831,1.828) (0.885,1.849) (0.940,1.869) (0.994,1.888) (1.049,1.907) (1.104,1.925) (1.160,1.943) (1.215,1.960) (1.271,1.976) (1.327,1.992) (1.383,2.007) (1.439,2.021) (1.495,2.034) (1.552,2.047) (1.609,2.059) (1.666,2.070) (1.723,2.080) (1.780,2.090) (1.837,2.099) (1.895,2.107) (1.952,2.114) (2.010,2.120) (2.068,2.126) (2.125,2.130) (2.183,2.134) (2.241,2.137) (2.299,2.139) (2.357,2.140) (2.415,2.140) (2.473,2.139) (2.531,2.137) (2.589,2.134) (2.647,2.131) (2.705,2.126) (2.763,2.120) (2.820,2.114) (2.878,2.106) (2.935,2.097) (2.992,2.088) (3.049,2.077) (3.106,2.065) (3.163,2.052) (3.219,2.039) (3.275,2.024) (3.331,2.008) (3.345,2.003)};
\draw[flinea, flecha=1.520cm] plot coordinates {(-0.091,1.495) (-0.143,1.521) (-0.195,1.546) (-0.247,1.572) (-0.299,1.597) (-0.352,1.622) (-0.404,1.646) (-0.457,1.670) (-0.510,1.694) (-0.563,1.718) (-0.616,1.741) (-0.670,1.763) (-0.723,1.786) (-0.777,1.807) (-0.831,1.828) (-0.885,1.849) (-0.940,1.869) (-0.994,1.888) (-1.049,1.907) (-1.104,1.925) (-1.160,1.943) (-1.215,1.960) (-1.271,1.976) (-1.327,1.992) (-1.383,2.007) (-1.439,2.021) (-1.495,2.034) (-1.552,2.047) (-1.609,2.059) (-1.666,2.070) (-1.723,2.080) (-1.780,2.090) (-1.837,2.099) (-1.895,2.107) (-1.952,2.114) (-2.010,2.120) (-2.068,2.126) (-2.125,2.130) (-2.183,2.134) (-2.241,2.137) (-2.299,2.139) (-2.357,2.140) (-2.415,2.140) (-2.473,2.139) (-2.531,2.137) (-2.589,2.134) (-2.647,2.131) (-2.705,2.126) (-2.763,2.120) (-2.820,2.114) (-2.878,2.106) (-2.935,2.097) (-2.992,2.088) (-3.049,2.077) (-3.106,2.065) (-3.163,2.052) (-3.219,2.039) (-3.275,2.024) (-3.331,2.008) (-3.345,2.003)};
\draw[flinea, flecha=1.657cm] plot coordinates {(0.099,1.473) (0.156,1.485) (0.212,1.498) (0.269,1.510) (0.326,1.522) (0.382,1.534) (0.439,1.545) (0.496,1.556) (0.553,1.567) (0.610,1.577) (0.668,1.586) (0.725,1.595) (0.782,1.603) (0.840,1.610) (0.898,1.617) (0.955,1.623) (1.013,1.628) (1.071,1.632) (1.129,1.635) (1.187,1.638) (1.245,1.639) (1.303,1.640) (1.361,1.640) (1.419,1.638) (1.477,1.636) (1.535,1.632) (1.592,1.628) (1.650,1.622) (1.708,1.615) (1.765,1.607) (1.822,1.598) (1.880,1.588) (1.936,1.576) (1.993,1.564) (2.049,1.550) (2.105,1.534) (2.161,1.517) (2.216,1.499) (2.270,1.480) (2.325,1.459) (2.378,1.437) (2.431,1.414) (2.484,1.389) (2.535,1.363) (2.586,1.335) (2.636,1.306) (2.686,1.275) (2.734,1.243) (2.781,1.210) (2.828,1.175) (2.873,1.138) (2.917,1.101) (2.960,1.061) (3.001,1.021) (3.041,0.979) (3.080,0.935) (3.116,0.891) (3.152,0.845) (3.185,0.797) (3.217,0.749) (3.247,0.699) (3.275,0.648) (3.301,0.597) (3.325,0.544) (3.336,0.517)};
\draw[flinea, flecha=1.657cm] plot coordinates {(-0.099,1.473) (-0.156,1.485) (-0.212,1.498) (-0.269,1.510) (-0.326,1.522) (-0.382,1.534) (-0.439,1.545) (-0.496,1.556) (-0.553,1.567) (-0.610,1.577) (-0.668,1.586) (-0.725,1.595) (-0.782,1.603) (-0.840,1.610) (-0.898,1.617) (-0.955,1.623) (-1.013,1.628) (-1.071,1.632) (-1.129,1.635) (-1.187,1.638) (-1.245,1.639) (-1.303,1.640) (-1.361,1.640) (-1.419,1.638) (-1.477,1.636) (-1.535,1.632) (-1.592,1.628) (-1.650,1.622) (-1.708,1.615) (-1.765,1.607) (-1.822,1.598) (-1.880,1.588) (-1.936,1.576) (-1.993,1.564) (-2.049,1.550) (-2.105,1.534) (-2.161,1.517) (-2.216,1.499) (-2.270,1.480) (-2.325,1.459) (-2.378,1.437) (-2.431,1.414) (-2.484,1.389) (-2.535,1.363) (-2.586,1.335) (-2.636,1.306) (-2.686,1.275) (-2.734,1.243) (-2.781,1.210) (-2.828,1.175) (-2.873,1.138) (-2.917,1.101) (-2.960,1.061) (-3.001,1.021) (-3.041,0.979) (-3.080,0.935) (-3.116,0.891) (-3.152,0.845) (-3.185,0.797) (-3.217,0.749) (-3.247,0.699) (-3.275,0.648) (-3.301,0.597) (-3.325,0.544) (-3.336,0.517)};
\draw[flinea, flecha=1.414cm] plot coordinates {(0.102,1.450) (0.159,1.450) (0.217,1.450) (0.275,1.449) (0.333,1.448) (0.391,1.447) (0.449,1.445) (0.507,1.443) (0.565,1.440) (0.623,1.436) (0.681,1.432) (0.739,1.427) (0.797,1.421) (0.854,1.415) (0.912,1.407) (0.969,1.399) (1.026,1.389) (1.083,1.379) (1.140,1.367) (1.197,1.354) (1.253,1.340) (1.309,1.325) (1.365,1.309) (1.420,1.291) (1.474,1.272) (1.529,1.251) (1.582,1.229) (1.635,1.205) (1.688,1.180) (1.739,1.154) (1.790,1.126) (1.840,1.096) (1.889,1.065) (1.937,1.032) (1.983,0.998) (2.029,0.961) (2.073,0.924) (2.115,0.884) (2.156,0.843) (2.196,0.801) (2.233,0.757) (2.269,0.711) (2.303,0.664) (2.334,0.615) (2.364,0.565) (2.391,0.514) (2.415,0.461) (2.437,0.408) (2.457,0.353) (2.473,0.298) (2.487,0.241) (2.498,0.184) (2.506,0.127) (2.511,0.069) (2.513,0.011) (2.514,-0.000)};
\draw[flinea, flecha=1.414cm] plot coordinates {(-0.102,1.450) (-0.159,1.450) (-0.217,1.450) (-0.275,1.449) (-0.333,1.448) (-0.391,1.447) (-0.449,1.445) (-0.507,1.443) (-0.565,1.440) (-0.623,1.436) (-0.681,1.432) (-0.739,1.427) (-0.797,1.421) (-0.854,1.415) (-0.912,1.407) (-0.969,1.399) (-1.026,1.389) (-1.083,1.379) (-1.140,1.367) (-1.197,1.354) (-1.253,1.340) (-1.309,1.325) (-1.365,1.309) (-1.420,1.291) (-1.474,1.272) (-1.529,1.251) (-1.582,1.229) (-1.635,1.205) (-1.688,1.180) (-1.739,1.154) (-1.790,1.126) (-1.840,1.096) (-1.889,1.065) (-1.937,1.032) (-1.983,0.998) (-2.029,0.961) (-2.073,0.924) (-2.115,0.884) (-2.156,0.843) (-2.196,0.801) (-2.233,0.757) (-2.269,0.711) (-2.303,0.664) (-2.334,0.615) (-2.364,0.565) (-2.391,0.514) (-2.415,0.461) (-2.437,0.408) (-2.457,0.353) (-2.473,0.298) (-2.487,0.241) (-2.498,0.184) (-2.506,0.127) (-2.511,0.069) (-2.513,0.011) (-2.514,-0.000)};
\draw[flinea, flecha=1.114cm] plot coordinates {(0.099,1.427) (0.155,1.414) (0.212,1.401) (0.268,1.388) (0.325,1.374) (0.381,1.360) (0.437,1.345) (0.493,1.330) (0.549,1.314) (0.604,1.297) (0.660,1.280) (0.715,1.262) (0.770,1.243) (0.824,1.223) (0.878,1.202) (0.932,1.179) (0.985,1.156) (1.037,1.131) (1.089,1.105) (1.140,1.078) (1.191,1.049) (1.240,1.019) (1.289,0.988) (1.336,0.955) (1.383,0.920) (1.428,0.884) (1.472,0.846) (1.514,0.806) (1.555,0.765) (1.594,0.722) (1.631,0.677) (1.666,0.631) (1.699,0.583) (1.730,0.534) (1.758,0.483) (1.784,0.431) (1.806,0.378) (1.826,0.324) (1.843,0.268) (1.857,0.212) (1.868,0.155) (1.875,0.097) (1.879,0.039) (1.880,0.000)};
\draw[flinea, flecha=1.114cm] plot coordinates {(-0.099,1.427) (-0.155,1.414) (-0.212,1.401) (-0.268,1.388) (-0.325,1.374) (-0.381,1.360) (-0.437,1.345) (-0.493,1.330) (-0.549,1.314) (-0.604,1.297) (-0.660,1.280) (-0.715,1.262) (-0.770,1.243) (-0.824,1.223) (-0.878,1.202) (-0.932,1.179) (-0.985,1.156) (-1.037,1.131) (-1.089,1.105) (-1.140,1.078) (-1.191,1.049) (-1.240,1.019) (-1.289,0.988) (-1.336,0.955) (-1.383,0.920) (-1.428,0.884) (-1.472,0.846) (-1.514,0.806) (-1.555,0.765) (-1.594,0.722) (-1.631,0.677) (-1.666,0.631) (-1.699,0.583) (-1.730,0.534) (-1.758,0.483) (-1.784,0.431) (-1.806,0.378) (-1.826,0.324) (-1.843,0.268) (-1.857,0.212) (-1.868,0.155) (-1.875,0.097) (-1.879,0.039) (-1.880,0.000)};
\draw[flinea, flecha=0.910cm] plot coordinates {(0.091,1.405) (0.143,1.379) (0.195,1.353) (0.246,1.327) (0.298,1.300) (0.349,1.273) (0.400,1.246) (0.451,1.218) (0.502,1.189) (0.552,1.160) (0.601,1.130) (0.650,1.099) (0.699,1.067) (0.747,1.034) (0.794,1.001) (0.840,0.966) (0.886,0.929) (0.930,0.892) (0.973,0.853) (1.015,0.813) (1.055,0.772) (1.094,0.729) (1.132,0.684) (1.167,0.638) (1.200,0.591) (1.232,0.542) (1.261,0.492) (1.287,0.440) (1.311,0.387) (1.331,0.333) (1.349,0.278) (1.364,0.222) (1.376,0.165) (1.384,0.108) (1.388,0.050) (1.390,0.000)};
\draw[flinea, flecha=0.910cm] plot coordinates {(-0.091,1.405) (-0.143,1.379) (-0.195,1.353) (-0.246,1.327) (-0.298,1.300) (-0.349,1.273) (-0.400,1.246) (-0.451,1.218) (-0.502,1.189) (-0.552,1.160) (-0.601,1.130) (-0.650,1.099) (-0.699,1.067) (-0.747,1.034) (-0.794,1.001) (-0.840,0.966) (-0.886,0.929) (-0.930,0.892) (-0.973,0.853) (-1.015,0.813) (-1.055,0.772) (-1.094,0.729) (-1.132,0.684) (-1.167,0.638) (-1.200,0.591) (-1.232,0.542) (-1.261,0.492) (-1.287,0.440) (-1.311,0.387) (-1.331,0.333) (-1.349,0.278) (-1.364,0.222) (-1.376,0.165) (-1.384,0.108) (-1.388,0.050) (-1.390,0.000)};
\draw[flinea, flecha=0.763cm] plot coordinates {(0.076,1.382) (0.119,1.344) (0.162,1.305) (0.205,1.266) (0.247,1.226) (0.290,1.187) (0.332,1.147) (0.374,1.107) (0.415,1.066) (0.455,1.024) (0.495,0.982) (0.535,0.940) (0.573,0.896) (0.611,0.852) (0.647,0.807) (0.682,0.761) (0.716,0.714) (0.748,0.665) (0.779,0.616) (0.808,0.566) (0.835,0.515) (0.860,0.462) (0.882,0.409) (0.902,0.354) (0.920,0.299) (0.934,0.243) (0.946,0.186) (0.955,0.129) (0.960,0.071) (0.962,0.013) (0.962,0.000)};
\draw[flinea, flecha=0.763cm] plot coordinates {(-0.076,1.382) (-0.119,1.344) (-0.162,1.305) (-0.205,1.266) (-0.247,1.226) (-0.290,1.187) (-0.332,1.147) (-0.374,1.107) (-0.415,1.066) (-0.455,1.024) (-0.495,0.982) (-0.535,0.940) (-0.573,0.896) (-0.611,0.852) (-0.647,0.807) (-0.682,0.761) (-0.716,0.714) (-0.748,0.665) (-0.779,0.616) (-0.808,0.566) (-0.835,0.515) (-0.860,0.462) (-0.882,0.409) (-0.902,0.354) (-0.920,0.299) (-0.934,0.243) (-0.946,0.186) (-0.955,0.129) (-0.960,0.071) (-0.962,0.013) (-0.962,0.000)};
\draw[flinea, flecha=0.652cm] plot coordinates {(0.046,1.360) (0.073,1.308) (0.099,1.257) (0.126,1.205) (0.152,1.153) (0.178,1.101) (0.203,1.049) (0.228,0.997) (0.253,0.944) (0.277,0.892) (0.301,0.839) (0.324,0.785) (0.346,0.732) (0.367,0.678) (0.387,0.623) (0.406,0.568) (0.423,0.513) (0.440,0.457) (0.454,0.401) (0.467,0.345) (0.479,0.288) (0.488,0.231) (0.495,0.173) (0.501,0.115) (0.504,0.058) (0.505,0.000)};
\draw[flinea, flecha=0.652cm] plot coordinates {(-0.046,1.360) (-0.073,1.308) (-0.099,1.257) (-0.126,1.205) (-0.152,1.153) (-0.178,1.101) (-0.203,1.049) (-0.228,0.997) (-0.253,0.944) (-0.277,0.892) (-0.301,0.839) (-0.324,0.785) (-0.346,0.732) (-0.367,0.678) (-0.387,0.623) (-0.406,0.568) (-0.423,0.513) (-0.440,0.457) (-0.454,0.401) (-0.467,0.345) (-0.479,0.288) (-0.488,0.231) (-0.495,0.173) (-0.501,0.115) (-0.504,0.058) (-0.505,0.000)};
\node[menos] at (2.514,-0.232) {};
\node[menos] at (-2.514,-0.232) {};
\node[menos] at (1.880,-0.232) {};
\node[menos] at (-1.880,-0.232) {};
\node[menos] at (1.390,-0.232) {};
\node[menos] at (-1.390,-0.232) {};
\node[menos] at (0.962,-0.232) {};
\node[menos] at (-0.962,-0.232) {};
\node[menos] at (0.505,-0.232) {};
\node[menos] at (-0.505,-0.232) {};

    \end{scope}
    \node[cargapos] at (0,\esc) {};
    \node[rotulo,anchor=south,fill=panel,inner sep=1pt,rounded corners=1pt] at (0,1.0*\esc+0.17) {$q$};
    % toma de tierra
    \draw[gristexto,line width=0.6pt] (-2.0*\esc,-0.45*\esc) -- (-2.0*\esc,-0.75*\esc);
    \tierra{(-2.0*\esc,-0.75*\esc)}
    \node[rotulogris,anchor=east] at (2.25*\esc,-0.6*\esc) {$\varphi=0$};
    \draw[azulborde,line width=0.5pt] (-2.3*\esc,-1.0*\esc) rectangle (2.3*\esc,2.75*\esc);
    \node[rotulogris,anchor=south west] at (-2.3*\esc,2.77*\esc)
      {(b) Solución: líneas $\perp$ al conductor};
  \end{scope}
\end{tikzpicture}
